\section{Printed inputs and their hypotheses}\label{app:inputs}
The tables specify the statements used with their printed parameters. A local version means the version proved in Proposition~\ref{prop:partii}, with every local coefficient restriction retained. Proof identities from the input are used only with their displayed original coefficients and supports. The new geometry, Euler domain, completed norm, row envelope, and fourth-moment parameter are proved in the body.

\small
\begin{longtable}{@{}>{\raggedright\arraybackslash}p{.20\textwidth}>{\raggedright\arraybackslash}p{.39\textwidth}>{\raggedright\arraybackslash}p{.35\textwidth}@{}}
\toprule
Input in \cite{openai2026q}&Hypotheses retained&Application and verification\\
\midrule
\endfirsthead
\toprule
Input in \cite{openai2026q}&Hypotheses retained&Application and verification\\
\midrule
\endhead
\bottomrule
\endfoot
Theorem 1.1; (2.1)&The whole primitive finite-order Hecke family over $\F$; poles excluded from $\beta^*$.&Gives $\beta^*\le7/8$ and the exact ceiling $a\le\beta^*$ for actual-zero bins.\\
Proposition 2.1&$1/2<\sigma<1$, $C(s)=s+c$, positive global gap, target-independent low and high losses, one $J,f$, bounded nonvanishing $H$ on $\Rea s>\sigma$.&Proposition~\ref{prop:continuation} states and proves this input. Section~\ref{sec:geometry} verifies its hypotheses, including algebraic endpoint equality.\\
\S6.1; (6.1)--(6.2)&Fixed ray group $T$; complete zero-extended presentations supported on fixed $S$; nonnegative nonzero annular $W_0,W_1$; arbitrary positive scales; $Q=q_{b_{\rm aux}}XY$; original masks.&The scales are positive powers of $Z$. Retaining the whole-index indicator linearly in the absolute series derives the marked separation in Theorem~\ref{thm:low}.\\
Lemma 4.5&Fixed logarithmic boxes; one common coefficient measure for the row Hilbert space; specified finite weighted moments; uniformly smooth profiles. Sobolev parameters enter profiles, not arithmetic coefficients.&The low proof keeps row, dual, and slot norms in one profile. Witness choices are smooth-parameter choices. Equations~\eqref{eq:heightallocation} and \eqref{eq:highquantified} separate external and internal orders.\\
Lemma 4.6&Fixed positive truncation width, Gaussian annular seminorms summable with exponential weights, fixed polynomial-scale range.&The joint Gaussian profile meets these conditions. Half-width truncation and a threshold control whole retained annuli.\\
Lemma 4.7&Schwartz Fourier kernel, fixed normalized radial derivatives; bounded real Mellin strips and fixed weighted derivatives of the tests.&Genuine raw tails are removed before off-coprime extension. Raising orders changes seminorms, not real exponents.\\
Lemma 4.8; (4.14)&Primitive nonprincipal finite-order character with conductor norm $Q$; $-1/10\le\Rea s\le11/10$. For principal characters use the regularized zeta function.&The small-row numerator is on $1/2$ with conductor $\ll_{\mathcal A}U$. Plain reflection uses its primitive functional equation and restores all redundant radicals.\\
Lemma 4.9&$a\in[1/2,1]$, $0<e<10^{-3}$; zero-free disk of radius $2-a-2e$ about $2+it$, or global $b\ge\beta^*$ and $\Rea s\ge b+v$, $v>0$.&Buffered disks are supplied by Lemma 8.1. Global reciprocals use $\beta^*+e$ or $\beta^*+20e$. The extended prime moment explicitly uses $b=(1+\kappa)/2\ge\beta^*$.\\
Lemma 4.10&Squarefree deletion radical, local coefficients of modulus at most one; fixed positive real lower bound for inverse deletions; bounded strips for upper deletions.&Conductors and deletion supports are tracked separately and are polynomial. The primitive conjugate numerator follows from its original presentation and inverse deletions.\\
Lemma 5.2&Nonzero element rows, fixed finite $S$-supported ray family, fixed unit and $S$-valuation sectors; good moving primes remain moving.&The completed norm keeps rows meeting $S$, fixes sectors, and uses $f=1$, puncture one.\\
Proposition 5.1; Lemmas 5.3--5.4; (5.27); (9.15)&Original completed cubic coefficients, fixed reflection modulus, genuine common kernel and original dual support; squarefree sieve indices and row-independent coefficients; original dual indices retain permitted $S$-primes.&Appendix~\ref{app:partii} uses them in local reflection and the hybrid norm. The hybrid cubic column may contain permitted $S$-primes, as \src{(14.6), p.96} explicitly allows.\\
Lemma 7.1; (7.6), (7.11)--(7.13), (7.17)&Nonzero sixth-power-free $u$, $(u,S)=1$, $\p\notin S$, $j=0,\ldots,5$, original unit phases and zero extensions.&All six cases are displayed in Section~\ref{sec:euler}. New bounds are derived from these exact formulas.\\
Lemma 7.1, (7.14)&$s_r\ge51/100$, $z_r\ge17/50$, $w_r\ge-1/100$, $s_r+w_r\ge1+\eps_0$, fixed $\eps_0>0$.&Moderate paths use $10e$, small paths use $1/3$, and large paths lie farther right. Lemma~\ref{lem:euler} replaces the second-domain bound.\\
(7.2)&Nonnegative nonzero annular $W_0$; $M(z)=\int_0^\infty\widetilde W_0(r)r^{z-1}\,dr$ is holomorphic for $\Rea z>0$, with $M(\sigma)>0$ for $0<\sigma<1$.&Section~\ref{sec:geometry} uses $M_0(1/6)>0$ and the nonnegative nonzero $W_1$, giving $c_S>0$ in \eqref{eq:scalar}.\\
\S10.1&Fixed $S$, the original physical mask $(u,S)=1$ and sixth-power-free rows.&Section~\ref{sec:geometry} uses that $u=1$ is the only principal numerator row. The other unit rows remain in the small-row range.\\
Lemmas 8.1--8.2; Proposition 8.3&The family $\mathcal X_u$, nonprincipal inducing rows, bounded $d,r,m$, $T_1\le U^{1/100}$, smooth annuli, common pure twists, one frequency allowance, and the stated height condition.&Section~\ref{sec:inputs} records the conclusions. Theorem~\ref{thm:count} checks the coefficient spans. Section~\ref{sec:geometry} chooses $\tau$ after internal orders, including the pre-saturation range $[0,22]$.\\
Lemma 10.2&Fixed compact real boxes; joint density with integrated and pointwise trace decay; arithmetic factor $Z^B(1+|\boldsymbol t|)^J$ with $B,J$ fixed before the tail order.&Equation~\eqref{eq:externaltail} uses the original external smooth kernels. Buffered axes stay restricted. Internal measures with only weighted $L^1$ control are not used for traces.\\
Lemma 11.1&$1/2<\sigma<1$; common positive pre-height saving; finite $A_\eta,B_\eta,\tau_{0,\eta}$ before $N$; $J,f$ independent of $T_1$.&Equation~\eqref{eq:highquantified} supplies this order. Choose $\tau$, then $N$, then the threshold.\\
Proposition 11.3&$1/2\le\sigma<1$ and the whole primitive Hecke family zero-free in $\Rea s>\sigma$.&Every boundary here is in this range. Finite deletions are nonzero for $\Rea s>0$; Section~\ref{sec:inputs} handles pole cancellation at one.\\
\end{longtable}

\begin{longtable}{@{}>{\raggedright\arraybackslash}p{.20\textwidth}>{\raggedright\arraybackslash}p{.39\textwidth}>{\raggedright\arraybackslash}p{.35\textwidth}@{}}
\toprule
Part II local input&Hypotheses retained&Why they hold here\\
\midrule
\endfirsthead
\toprule
Part II local input&Hypotheses retained&Why they hold here\\
\midrule
\endhead
\bottomrule
\endfoot
Lemma 13.1&Fixed ray quotient, annular test, positive scale tending to infinity, fixed finite deletions.&Each $P_i=Z^{\ell_i}$ has $\ell_i>0$ and its nonnegative bump is fixed before the target.\\
Lemmas 13.2--13.4&Complete finite moduli, original principal zero convention, all common factors extracted; CRT on genuinely coprime residual products.&The transforms keep full products, shared-prime multiplicities, roots, masks, and kernels before separation.\\
Corollary 14.1; Lemmas 14.2--14.3&Fixed finite-ray coefficients and row-independent puncture; disjoint product-form lists; squarefree hybrid indices; bounded lengths; row restrictions independent of active slots and duals apart from the exact zero mask.&Theorem~\ref{thm:low} displays the actual coefficient and joint profile. The frozen rescaled tuple determines the row ball. The canonical terminal retains its allowed puncture and label.\\
Proposition 15.2&Original additive ray coefficients, fixed ball, polynomial $Q,Y'\ge1$, and $P_a=Y'^2/Q\ge1$.&The low proof gives the actual scales. Its inequality $l_y-\ell-11b/6>0$ bounds the last Gram term by $P_a^{1/6}$.\\
Proposition 16.1&$51/100\le a\le1$, central contour, nonprincipal numerator, primitive conjugate bound \eqref{eq:correctnumerator}, nonzero central selected factors.&The conjugation-closed family and common buffer give the correct primitive bound. Lemma~\ref{lem:cutoff} supplies the cutoff. Apply the strict-prime conductor saving once.\\
Lemma 17.1&Fixed finite-ray presentation, either common orientation, bounded row-independent prime coefficients, disjoint supports; $r+2z\le m-c_1$, $2r+8z\le3m-c_2$, fixed $c_1,c_2>0$.&At width one, Theorem~\ref{thm:count} gives $c_1\ge2\nu_0$ and $c_2\ge8\nu_0+9/37$. Whole conjugations preserve all zeros.\\
Lemma 17.2&Fixed character and puncture independent of row and label; squarefree divisor-bounded label weight depending only on that label; product-form mark; two strict puncture and width conditions.&Appendix~\ref{app:partii} states these conditions and the initialization margins. No arbitrary extra column coefficient is introduced.\\
Lemma 17.3; Corollary 17.4&Finite acyclic graph, common coefficient measure, fixed exponents and height maps, finite weighted moments.&Both inductions strictly decrease width and retain common full-product profiles. Internal orders are chosen backwards.\\
Lemma 17.5&Primitive finite character with zero extension, complete mask, radial Schwartz profile, positive row scale; principal restoration on nonempty scale at least one.&The canonical transforms use genuine masks and untruncated divisor sums. Restored principal frequencies keep the outer mask.\\
Lemma 17.6&Sixth-power-free rows, original inverse coefficients and zeros; $U,D\ge1$, $H=\max(2U,D^{1+c})$, $c>0$; bounded real profile powers and rowwise twists.&These are the physical rows. The long and zero-capacity branches use small fixed $c$, including through $r=1$.\\
Lemma 18.1, printed range&Same row and mask in both plains and slots; $\Theta$ coefficient span; inducing character outside $\Theta$ for positive slots; fixed-count-independent mesh; $n_1+n_2+6\kappa z\le M$ and the stated zero-free premise.&The original range is used at $\kappa=3/4$, or at $1$ for zero slots. Theorem~\ref{thm:fourth} proves the extension. The cubic stage invokes it with $\kappa_a=2\beta^*-1$.\\
Lemma 18.2&Full common coefficient, complete Möbius allocations, natural zeros and whole-product phases, equal-product formal rectangles.&Every extraction in Section~\ref{sec:fourth} retains that class and does not repuncture divisor quotients.\\
Lemma 18.3&Fixed character, common polynomial-norm squarefree mask, same norm power within a side, annular profiles, equal-product formal scales; all-scale fallback when lower scales fail.&Exceptional children stay formal rather than clipped and have the same coefficients and masks in both rectangles.\\
Lemma 19.1&Lemma 8.2's buffer and heights; prime-annulus frequency in the same allowance; original identity-ray slots.&Section~\ref{sec:count} uses that class. Amplitude subdivision follows the original-bin contour move.\\
(12.5), (16.2)--(16.4), (16.7)&Distinct finite slot tuples, original positive windows and masks, arbitrary positive scales; marked coefficient $\overline{\eta(\p)}$.&Equations~\eqref{eq:physical} and \eqref{eq:localidentity} identify the same finite probe. Full quotient-free continuation is proved here.\\
Lemma 20.2, algebra only&Original numerical function, $0\le\delta\le5/6$, $0\le x\le1/2$, positive $J$.&Equations~\eqref{eq:oldpoly}--\eqref{eq:oldsquare} reproduce its identity and signs. The changed geometry has an explicit increment.\\
\end{longtable}
\normalsize
